\section*{ID10211: Landauer Principle}
\paragraph{Metadata.}
PiID: 5320186837937 | RTEID: RTE:P39651.5806:T2.5018 | UUID: 477800a4-78f0-581b-a8c9-01c2ea22992d

\paragraph{Canonical Statement.}
to the environment at temperature \$T\$. The total power consumption for information processing is bounded by: \textbackslash{}begin\{equation\} P\_\{\textbackslash{}text\{info\}\} \textbackslash{}geq \textbackslash{}frac\{k\_B T \textbackslash{}ln 2\}\{\textbackslash{}tau\_\{\textbackslash{}text\{bit\}\}\} N\_\{\textbackslash{}text\{bits\}\} \textbackslash{}end\{equation\}

\paragraph{Canonical Equation.}
\[
P_{\text{info}} \geq \frac{k_B T \ln 2}{\tau_{\text{bit}}} N_{\text{bits}}
\]

\paragraph{Dictionary.}
Landauer Principle — P\_info ≥ frack\_B T ln 2τ\_bit N\_bits

\paragraph{Encyclopedia.}
Landauer Principle is an inequality / bound, written P\_info ≥ frack\_B T ln 2τ\_bit N\_bits. It states an inequality or bound that constrains the admissible range of the quantity. It is built from P\_, k\_B, T, N\_. Within the Trawin Zero Point registry it serves as one element of the framework's mathematical narrative.
